\documentclass[10pt,a4paper]{amsart}
\usepackage[margin=19mm]{geometry}
\usepackage{amsmath,amssymb,mathtools}
\usepackage[scr=rsfs,cal=euler]{mathalfa}
\usepackage{xfrac}
\usepackage[unicode,hidelinks,bookmarks=false]{hyperref}
\newcommand{\ie}{i.e.\ }
\newcommand{\cf}{cf.\ }
\newcommand{\resp}{respectively\ }
\newcommand{\Subsection}{\subsection}
\providecommand{\nobreakdash}{\nobreak\hbox{-}}
\setlength{\emergencystretch}{2em}
\multlinegap0pt
\usepackage{amssymb}
\usepackage{tikz}
\usepackage{enumerate}
\usepackage[shortlabels]{enumitem}
\usepackage{float}
\newtheorem{lemma}{Lemma}
\newcommand{\N}{{\mathbb N}}
\newcommand{\R}{{\mathbb R}}       % Field of real numbers
\newcommand{\Sp}{{\mathbb S}}
\newcommand{\Z}{{\mathbb Z}}    
\title{An alternate approach to bilinear rough singular integrals}
\author{Ankit Bhojak, Saurabh Shrivastava}
\date{Source: arXiv:2508.19181v3 (CC BY 4.0)}
\begin{document}
\maketitle

\noindent\emph{Source and licence.} An alternate approach to bilinear rough singular integrals. Version arXiv:2508.19181v3 (CC BY 4.0), distributed by arXiv under the Creative Commons Attribution 4.0 International licence, http://creativecommons.org/licenses/by/4.0/, as recorded in the arXiv licence metadata for this exact version and shown on the abstract page https://arxiv.org/abs/2508.19181v3. Full licence notice: \texttt{licenses/arxiv-2508.19181-CC-BY-4.0.txt}.\par\smallskip

\noindent\textit{Editorial scope.} Counted unit: the article's lemma lemma:ellinfty (source0.tex lines 475--478). The article's complete original proof of it is delivered with it (source0.tex lines 479--501, one \texttt{proof} environment of 23 lines). No mathematics is written, repaired or re-derived: the statement and the proof are the article's own lines, in the article's own notation, and no retained line is substituted or re-flowed.  No further source-local context is needed: every cross-reference of every retained line resolves inside the two delivered spans.  Nothing was omitted from any retained span, so the package carries no omission record.  No retained line contains a citation, so the wrapper carries no bibliography and no bibliography entry.  No retained line contains a figure environment, an image-inclusion command or drawing code, so no image is delivered and the package declares no asset dependency.  Editorial formatting only, in addition: an AMS \texttt{amsart} preamble that restates the article's own declarations and macros used by the retained lines, namely the environments \texttt{lemma} on separate counters, together with the standard packages those lines need.  The retained environments print in this document as Lemma 1 in order of appearance, whereas the article prints the same items with the numbering of its own full document.  The complete native source of the article as distributed in the arXiv e-print for this exact version is bundled unchanged as \texttt{sources/183937/source0.tex}.
\begin{lemma}\label{lemma:ellinfty}
	For any $0<\epsilon_2<1$ and $k_1,k_2,k_3\in\Z^d$ with $|(k_1,k_2)|\geq\frac{3}{8}\lambda$, we have
	\[|I_{\vec{\mathbf{k}}}|\lesssim\lambda^{-(1-\epsilon_2)}\|\Omega\|_\infty.\]
\end{lemma}

\begin{proof}
	For a fixed $0<\epsilon_2<1$, we write $I_{\vec{\mathbf{k}}}=I_{\vec{\mathbf{k}}}^++I_{\vec{\mathbf{k}}}^-$, where
\begin{align*}
	I_{\vec{\mathbf{k}}}^+&=\int_{\{(\theta_1,\theta_2)\in\Sp^{2d-1}:|k_1.\theta_1+k_2.\theta_2|>\lambda^{\epsilon_2}\}}\Omega(\theta)\int_{\R^d}\int_{\frac{1}{2}}^2\zeta_{\theta}(x,r)e^{-2\pi i\mathcal{P}_{\vec{\mathbf{k}},\theta}(x,r)}drdxd\sigma(\theta),\\
	I_{\vec{\mathbf{k}}}^-&=I_{\vec{\mathbf{k}}}-I_{\vec{\mathbf{k}}}^+.
\end{align*}
For the term $I_{\vec{\mathbf{k}}}^+$, we apply $N-$fold integration by parts in the $r-$variable to obtain
\begin{align*}
	\left|I_{\vec{\mathbf{k}}}^+\right|&\lesssim \int_{\{|k_1.\theta_1+k_2.\theta_2|>\lambda^{\epsilon_2}\}}|\Omega(\theta)|\int_{\R^d}\int_{\frac{1}{2}}^2\frac{|\partial_{r}^N\zeta_{\theta}(x,r)|}{|k_1.\theta_1+k_2.\theta_2|^N}drdxd\sigma(\theta)\lesssim \frac{1}{\lambda^{\epsilon_2 N}}\|\Omega\|_1.
\end{align*}
We choose $N\in\N$ such that $\epsilon_2 N>100$.
We observe that the following elementary estimate holds. 
\begin{equation*}
	N(\omega,\varepsilon):=\sigma(\{\theta\in\Sp^{2d-1}:|\omega\cdot\theta|\leq\varepsilon\})\lesssim\varepsilon,\quad \text{for}\;\omega\in\Sp^{2d-1}.
\end{equation*}
Further, note that for $\frac{3}{8}\lambda\leq|(k_1,k_2)|$ and $|k_1\theta_1+k_2\theta_2|\leq\lambda^{\epsilon_2}$, we have that \[\left|\left(\frac{(k_1,k_2)}{|(k_1,k_2)|}\cdot(\theta_1,\theta_2)\right)\right|\lesssim\lambda^{-1+\epsilon_2}.\] 
In view of the above estimates, we obtain that
\begin{align*}
	|I_{\vec{\mathbf{k}}}^-|\lesssim\;&\int\limits_{\substack{\theta\in\Sp^{2d-1}:\\\left|\left(\frac{(k_1,k_2)}{|(k_1,k_2)|}\cdot(\theta_1,\theta_2)\right)\right|\lesssim\lambda^{-1+\epsilon_2}}}|\Omega(\theta)|\int_{\R^d}\int_{\frac{1}{2}}^2|\zeta_{\theta}(x,r)|drdxd\sigma(\theta)\\
	\lesssim\;&\|\Omega\|_\infty N\left(\frac{(k_1,k_2)}{|(k_1,k_2)|},4\lambda^{-1+\epsilon_2}\right)\\
	\lesssim\;&\lambda^{-(1-\epsilon_2)}\|\Omega\|_\infty.
\end{align*}
\end{proof}

\end{document}
